
\documentclass[a4paper,french]{rnti}


\usepackage[T1]{fontenc}
\usepackage[utf8]{inputenc}

%pour bien présenter les URL et autres adresses emails
\usepackage{url}

\usepackage{graphicx}

\usepackage{hyperref}

% Titre court pour entête
\titrecourt{Regards d'actualité sur les données historisées d'EGC}

% Noms auteurs pour entête :
%    Si un seul auteur, mettre : Initiale. NomPremierAuteur
%    Si deux auteurs, mettre : Initiale1. NomPremierAuteur et Initiale1. 
NomDeuxiemeAuteur
%    Si plus de deux auteurs, mettre comme ci-dessous
%
\nomcourt{S. Valat et C. Favre}


\titre{Regards d'actualité au prisme des enjeux sociétaux \\sur les 
données historisées d'EGC
}

\auteur{Sébastien Valat\affil{1},
        Cécile Favre\affil{2}\affilsep\affil{3}\\
        }

\affiliation{
    \affil{1}ATOS\\
          sebastien.valat@atos.net\\
          \http{https://atos.net/fr/}\\
    %
    \affil{2}Université de Lyon, Lyon 2, ERIC EA 3083 (laboratoire 
d'informatique)\\
          cecile.favre@univ-lyon2.fr\\
          \http{https://eric.msh-lse.fr/}\\
    %
    \affil{3}Université de Lyon, Lyon 2, CMW UMR 5283 (laboratoire de 
sociologie)\\
                  \http{https://www.centre-max-weber.fr/}\\
                  
     
          
 }

\resume{
Au-delà du contenu scientifique produit dans le cadre d'EGC, EGC désigne 
aussi une communauté scientifique composée de personnes qui s'inscrivent 
dans un environnement traversé par des enjeux sociétaux qui interpellent 
en tant que citoyen·nes. Partant de cette réalité et de la demande du défi 
« 20 ans d’histoire pour quel avenir ? », il a paru pertinent de nous 
focaliser sur comment EGC pouvait se saisir de certains de ces enjeux. 
Nous axons notre analyse au regard de trois enjeux : le rapport au 
travail, l’égalité femmes-hommes et l’écologie, invitant ainsi à une 
réflexion plus large sur nos postures professionnelles.

}

\summary{
Beyond the scientific content produced in the framework of the EGC 
conference, EGC is also a scientific community composed of people who are 
part of a social environment crossed by different societal issues, that 
challenge as citizens. Starting from this reality, as well as from the 
challenge of "20 years of history for which future?", it seemed 
appropriate to focus on how EGC could address some of these issues. We 
focus our analysis on three issues: the relationship to work, gender 
equality and ecology, inviting us to a broader reflection on our 
professional postures.
%
}


\begin{document}


\section{Introduction}\label{sec: introduction}

La conférence EGC a été créée en 2001 et 19 ans ont passé. Elle a su 
fédérer autour de son existence une communauté scientifique qui rassemble 
des personnes dont les intérêts de recherche la définissent. Cette 
communauté est constituée de personnes qui contribuent scientifiquement à 
la conférence et qui sont par ailleurs des citoyen·nes de la société, 
s'insérant dans un environnement social traversé par différents enjeux 
sociétaux pouvant nous interpeller.



Comme le précisait \cite{pestre} dans l'\textit{Introduction aux Sciences 
studies} : « Malgré ce qu'énonce la philosophie spontanée des savants, il 
est difficile de laisser la science dans sa tour d'ivoire et de faire 
comme si c'était dans l'univers clos des laboratoires que se passait 
l'essentiel de ce qui la concerne. [...], la question du rapport des 
sciences [...] à l'économique et au politique (comme aux sociabilités et à 
la domination masculine) ne peut être évacuée ».


 Partant de là, et de la demande du défi « 20 ans d’histoire pour quel 
avenir ? », nous avons souhaité adopter une perspective particulière pour 
nous focaliser sur comment EGC pouvait se saisir de certains de ces enjeux 
sociétaux. En effet, dans cet article, nous ne concentrons pas l’analyse 
sur le contenu scientifique produit (par exemple en termes d'analyse de 
l'évolution de thématiques, ou de la dimension collaborative), mais nous 
faisons le choix d’axer notre analyse au regard de trois enjeux sociétaux 
qui nous paraissent fondamentaux : le rapport au travail (articulation des 
temps de vie), l’égalité femmes-hommes et l’écologie.

Les travaux scientifiques en sciences humaines et sociales se multiplient 
autour de l'épuisement au travail, de l'articulation des temps de vie, 
sans que le métier d'enseignant·e-chercheur·e ne soit épargné par ces 
problématiques, et ce dans différents pays (\citep{souffrance-ec-bresil15, 
souffrance-ec-tunisie18}. L'usage du mail apparaît être comme une des 
composantes pertinentes à analyser. C'est ce que font d'ailleurs 
\cite{mail-ec14} dans un article intitulé \textit{Moduler sa connexion : 
les enseignants-chercheurs aux prises avec leur courriel}, analysant sur 
la base d'entretiens l'usage du mail. Au vu des données fournies pour le 
défi EGC sur les mails de la liste de diffusion, il nous a paru pertinent 
de faire un focus sur l’analyse des heures d’envoi de ces mails, pour 
traiter de cette question du temps de travail.

Par ailleurs, l’enjeu de l’égalité femmes-hommes (déclarée grande cause du 
quinquennat du président de la république actuel), s’inscrit dans une 
société française qui a vu une médiatisation croissante sur le sujet. 
\cite{collet19book} livre dans son dernier ouvrage intitulé \textit{Les 
oubliées du numérique} une analyse de la place des femmes dans le 
numérique, de façon historisée notamment. Si les efforts menés pour 
attirer davantage de femmes dans le domaine de l’informatique sont de plus 
en plus importants, l’enjeu de leur visibilité est également présent (en 
termes d'égalité professionnelle, en particulier pour l'évolution de 
carrière par rapport au plafond de verre). Nous nous intéressons alors à 
la place de celles-ci dans la communauté EGC. Ceci nous a amené·es à 
étiqueter manuellement par rapport au sexe non seulement les données 
fournies sur les auteur/trices, mais aussi des données complémentaires que 
nous avons récoltées sur les conférences invitées, les comités de 
programme et les différentes présidences.


Enfin, la dimension écologique constitue aujourd’hui un enjeu fondamental, 
notamment du point de vue du réchauffement climatique. Alors même que la 
réflexion globale en matière d’écologie numérique est croissante, il 
s’agit d’amorcer des réflexions vis-à-vis des émissions de carbone, et le 
transport constitue un des domaines où des actions sont et doivent être 
entreprises. L'ADEME (Agence De l’Environnement et de la Maîtrise de 
l’Energie), établissement public sous la tutelle conjointe du ministère de 
la Transition écologique et solidaire et du ministère de l’Enseignement 
supérieur, de la Recherche et de l’Innovation, précise qu'en France il est 
considéré que les transports sont responsables d’1/3 des émissions de gaz 
à effet de serre. Dans ce papier, nous nous intéressons alors à la 
dimension spatiale du point de vue du déplacement des auteur/trices aux 
différentes éditions de la conférence EGC. 

Cet article est alors organisé de la façon suivante. Dans la section , 
nous précisons le travail de récupération et de préparation des données de 
manière générale. Puis, dans la section \ref{sec: rapport_travail}, nous 
abordons la dimension travail avec l'analyse des heures d'envoi de mail. 
Ensuite, dans la section \ref{sec:egalite}, nous analysons la place des 
femmes dans la communauté EGC. Nous proposons ensuite, dans la section , 
d'évoquer la dimension écologique avec l'analyse en termes de transport 
des auteur/trices. Enfin nous apportons une conclusion de ce travail dans 
la section .

\section{Récupération et préparation des 
données}\label{sec:recup-prepa-data}

Notre objectif était de pouvoir exploiter le maximum de données pour 
traiter les axes choisis, allant au-delà parfois des données fournies. La 
période considérée pour les trois thématiques n'est pas systématiquement 
la même en fonction des données accessibles.


Pour la dimension travail, les mails fournis vont de juillet 2006 à 
septembre 2018.

Sur la dimension de la place des femmes, nous avons cherché manuellement, 
pour compléter les données fournies, des données sur le site Web d'EGC et 
les sites Web des éditions. Nous avons récolté ainsi les personnes ayant 
présenté les conférences invitées, la composition des comités de programme 
et les différentes présidences : les personnes invitées à la présidence 
d'honneur et celles qui ont présidé les comités de programme et 
d'organisation, nous permettant de faire une analyse de 2001 à 2019 
inclus. Un important travail a été réalisé pour procéder à un étiquetage 
manuel de la variable sexe. L'objectif était d'avoir des données précises, 
là où les algorithmes d'étiquetage automatique dédiés ne produisent pas un 
résultat à l'unité près.




Concernant la dimension écologique, celle-ci a nécessité un travail 
important de préparation en développant un processus le plus automatisé 
possible. Il s'agissait de partir des documents PDF eux-mêmes pour accéder 
à l'affiliation des auteur/trices.
Ainsi, compte-tenu de la disponibilité des données sur RNTI, cette analyse 
démarre à partir de 2006 et va jusqu'en 2019. La préparation de ces 
données nécessitant davantage de développements informatiques, des 
éléments seront précisés dans la partie dédiée à cette thématique en 
section \ref{sec: ecologie}.



\section{Focus rapport au travail}\label{sec: rapport_travail}




Le droit à la déconnexion est une disposition issue de la loi El Khomri de 
2016, dite « loi Travail ». Elle est entrée en vigueur le 1er janvier 
2017. Un des objectifs de cette loi est d'adapter le droit du travail à 
l'ère du numérique.  Concernant en particulier le droit à la déconnexion, 
l'objectif est de permettre aux salarié·es de concilier vie personnelle et 
vie professionnelle, tout en luttant contre les risques de burn out. Pour 
cela, ils/elles doivent avoir la possibilité de ne pas se connecter aux 
outils numériques et de ne pas être contacté·es par leur employeur en 
dehors de leur temps de travail (congés payés, jours de RTT, week-end, 
soirées...). 

Ce droit à la déconnexion concerne toutes les personnes salariées. 
Cependant, il est à noter que cette loi ne s'applique pas dans la fonction 
publique de façon réglementaire. Des initiatives peuvent néanmoins être 
prises dans les institutions. La question se pose alors sur la capacité à 
respecter ce temps de déconnexion, dans des environnements dont le temps 
de travail peut être une variable assez « élastique », et c'est le cas 
pour beaucoup d'enseignant·es-chercheur·es.






Le jeu de données de la liste de diffusion EGC comporte 7075 mails écrits 
entre le lundi 10 juillet 2006 et le dimanche 9 septembre 2018, provenant 
de 811 adresses mails différentes. 




Pour travailler sur les horaires, nous avons utilisé le champ 
\textit{date} des emails considérant qu'il reflète normalement l'heure 
locale\footnote{en tenant compte du fuseau horaire local selon la norme 
RFC 4021, sous réserve que les mails envoyés ne l'aient pas été durant un 
déplacement à l'étranger avec décalage horaire} pour la personne qui 
envoie le mail, contrairement au champ \textit{received} donnant la date 
du côté serveur SMTP potentiellement distant.

Nous avons choisi de considérer l'envoi de mails dans le cadre 
professionnel de 8h à 20h comme « classique ». Nous avons pu dénombrer 
1116 mails écrits en dehors de cette plage horaire, soit 16\% des mails. 
Ceci est une trace du travail non négligeable effectué en dehors 
d'horaires classiques de travail, car si les mails aux listes de diffusion 
peuvent être considérés marginaux dans l'activité professionnelle, ces 
derniers sont la trace de cette activité.


 

Dans la figure \ref{fig: mails_out} à la page \pageref{fig: mails_out}, 
nous pouvons constater que l'augmentation de l'activité sur la liste de 
diffusion va de paire avec une augmentation du pourcentage de mails 
envoyés hors horaires classiques, qui oscille à présent autour de 15\%.

\begin{figure}[h]\centering
\includegraphics[width=0.65\linewidth,keepaspectratio]{images/mails_out.png}
\caption{\'{E}volution des mails postés sur la liste de diffusion EGC en 
fonction des plages horaires dédiées ou non classiquement au 
travail.}\label{fig: mails_out}
\end{figure}

Pour aller plus en détails, nous avons découpé la période non classique en 
3 créneaux : de 20h à minuit, de minuit à 4h et de 4h à 8h. Les résultats 
sont fournis dans la figure~\ref{fig: mails_out}. Nous pouvons y observer 
que la tranche 20h-minuit est celle qui est la plus utilisée, 
correspondant à une activité de travail en soirée. Parmi les 1116 mails 
hors créneaux classiques, 70\% ont été écrits de 20h à minuit ; 15\% de 
minuit à 4h et 15\% de 4h à 8h. Ce travail en dehors d'horaires classiques 
doit nous interpeller sur le rapport au travail et notre organisation du 
temps de travail.

\begin{figure}[h]\centering
\includegraphics[width=0.65\linewidth,keepaspectratio]{images/mails_creneaux_out.png}
\caption{Evolution des mails postés sur la liste de diffusion EGC en 
fonction des plages horaires et des années pour les mails postés en dehors 
des plages classiques de travail.}\label{fig:mails_creneaux_out}
\end{figure}



\section{Focus égalité femmes-hommes}\label{sec:egalite}


 
La place des femmes dans l'enseignement supérieur et la 
recherche\footnote{Rédaction en 2014 d’un Livre Blanc « Le genre dans 
l’enseignement supérieur et la recherche » par des chercheur·es de l’ANEF 
(Association nationale des études féministes)} reste une thématique 
d'actualité. Lors du défi 2016 sur les données d'EGC jusqu'à 2015, un 
premier travail avait été réalisé autour des femmes dans la communauté EGC 
\citep{cabanac16egc}. Il s'agit ici de reprendre des éléments d'analyse en 
prenant en compte les années depuis ce précédent défi et porter notre 
attention sur des aspects qui n'avaient pas été traités précédemment.


Tout d'abord, pour recontextualiser l'analyse, il s'agit de préciser que 
la communauté EGC concerne des personnes qui seraient plutôt catégorisées 
en sections 27 et 26 en terme de nomenclature CNU.
D'après les dernières fiches démographiques disponibles, soit pour 
2017-2018, fournies par le 
ministère\footnote{https://www.enseignementsup-recherche.gouv.fr/cid85019/fiches-demographiques-des-sections-du-cnu.html}, 
la section 27 comprenait 24\% de femmes et la section 26 comprenait 27\% 
de femmes (PR et MCF confondus). Ceci peut servir de base de comparaison 
sur la composition actuelle en matière de représentation sexuée, 
positionnant ainsi les résultats en perspective du « vivier » 
d'enseignant·es-chercheur·es, même si les statuts de chercheur·es (CR et 
DR) ne sont pas considérés dans ces données du ministère.



\subsection{Les auteurs et autrices d'EGC}\label{sec: auteurices}

Sur 19 éditions, parmi les co-auteur/trices des papiers (toutes catégories 
confondues qui sont présents dans les actes),  nous trouvons 27\% de 
femmes en moyenne.



Dans la figure \ref{fig: nb_auteurs_autrices}, nous pouvons observer 
l'évolution de cette répartition sexuée en nombre et en pourcentage. Le 
pourcentage a oscillé entre 21 et 31\%. Ainsi, nous pouvons considérer que 
la communauté EGC est dans la moyenne de ce qui s'observe en termes de 
répartition sexuée des sections CNU concernées.

\begin{figure}[h]\centering
\includegraphics[width=0.65\linewidth,keepaspectratio]{images/nb_auteurs_autrices.png}
\caption{Représentation sexuée des auteur/trices sur la période 
2001-2019.}\label{fig: nb_auteurs_autrices}
\end{figure}




Considérons à présent les personnes qui ont le plus contribué à la 
conférence sur les 19 ans, en nombre de papiers (quelle que soit la taille 
des papiers et moyennant les choix d'édition de certaines 
années\footnote{Par exemple, lors de certaines éditions, les actes n'ont 
pas comporté les papiers « démonstration de logiciel »}). La personne qui 
a le plus de papiers dans EGC sur la période 2001-2019 est un homme. Son 
premier papier remonte à la 2ème édition d'EGC, et il a en tout publié 37 
papiers (hors papier extension pour la version étendue à l'international).

Dans la figure \ref{fig: nb_contrib_max} de la page \pageref{fig: 
nb_contrib_max}, nous pouvons observer la représentation sexuée  des 
personnes ayant le plus contribué à EGC. Si l'on prend en compte les 
personnes ayant au moins 10 papiers sur les 19 éditions, elles sont 48, 
dont 12 femmes (soit 25\%). Nous restons donc sur les mêmes ordres de 
grandeur, ce qui montre que les femmes font bien partie aussi des 
personnes contribuant de façon importante et régulière à EGC.

\begin{figure}[h]\centering
\includegraphics[width=0.65\linewidth,keepaspectratio]{images/nb_contrib_max.png}
\caption{Représentation sexuée des auteur/trices ayant le plus contribué 
sur 2001-2019.}\label{fig: nb_contrib_max}
\end{figure}




\subsection{Présidences} \label{sec: presidences}

Un autre aspect intéressant est le travail et la visibilité qui découlent 
du rôle dans une présidence. Nous avons recensé les présidences d'honneur, 
de comité de programme (CP) et d'organisation (CO). Dans la figure 
\ref{fig: presidences}, nous pouvons observer le nombre de femmes et 
d'hommes impliqué·es. Ainsi, le pourcentage de femmes est de 17\% pour la 
présidence d'honneur, 24\% pour la présidence de comité de programme et 
33\% pour la présidence de comité d'organisation. Il est à noter que, 
parfois, des binômes de personnes ont pu être constitués, soit pour la 
gestion du comité de programme (surtout dans les premières éditions), soit 
pour la gestion de l'organisation. Si l'on considère le « prestige » 
pouvant être associé aux différents types de présidence, la dimension 
scientifique pouvant être globalement plus valorisée, il est intéressant 
de noter que l'on retrouve davantage de femmes sur l'aspect 
organisationnel des conférences, qui pourrait être mis en perspective 
d'analyses sociologiques sur la question des tâches et du genre.




\begin{figure}[h]\centering
\includegraphics[width=0.65\linewidth,keepaspectratio]{images/presidences.png}
\caption{Représentation sexuée pour la présidence d'honneur, la présidence 
du Comité de Programme et la présidence du Comité 
d'Organisation.}\label{fig: presidences}
\end{figure}




\subsection{Membres du comité de programme}\label{sec: comite}

Nous nous sommes également penché·es sur la présence des femmes dans le 
comité de programme. En étude sur le genre, un des indicateurs 
quantitatifs utilisés est le rapport de masculinité. Il est exprimé en 
nombre d'hommes pour 100 femmes (à la naissance, il est classiquement de 
105 garçons pour 100 filles). Dans la figure~\ref{fig:comite-programme}, 
nous constatons que, depuis 2015, ce rapport de masculinité a diminué 
jusqu'à 200 hommes pour 100 femmes, tendant vers un meilleur équilibre 
sexué sur la place des femmes dans le comité de programme. En 2019, cela 
représente 34\% de femmes dans le CP, c'est une manière de reconnaître 
leur compétence.

 

\begin{figure}[h]\centering
\includegraphics[width=0.65\linewidth,keepaspectratio]{images/CP.png}
\caption{Composition sexuée du Comité de Programme par année et rapport de 
masculinité.}\label{fig:CP}
\end{figure}



\subsection{Où sont les conférencières invitées ?}\label{sec: 
conferencieres}

Dans les carrières académiques, la visibilité / le rayonnement 
scientifique sont des éléments importants, notamment du point de vue de 
l'avancement dans la carrière, ou l'accès à des primes telles que la PEDR 
(Prime d'Encadrement Doctoral et de Recherche).

Il apparaît que l'invitation pour présenter une conférence invitée dans le 
cadre d'un colloque peut contribuer à cet enjeu de visibilité. Ainsi, bien 
que cet aspect ne figure pas dans les données initiales fournies pour le 
défi, nous avons choisi de les collecter manuellement à la fois via le 
site Web EGC qui récapitule partiellement ces éléments dans le menu dédié 
\footnote{https://www.egc.asso.fr/category/publications/conferences-invitees}, 
et nous avons complété les données qui s'avéraient lacunaires par une 
recherche dans les publications, compte-tenu du fait que les conférences 
invitées font l'objet d'un résumé d'une page maximum dans les actes de la 
conférence. Pour l'année 2001 et 2002, nous n'avons pu être en mesure de 
déterminer l'éventuelle présence de personnes donnant une conférence 
invitée, supposant que sur les deux premières éditions, il était possible 
qu'il n'y en ait pas eu.

Ainsi, entre 2003 et 2019, 65 personnes se sont succédées en session 
plénière pour présenter une conférence invitée, selon la répartition par 
année présentée dans la figure~\ref{sec:conf_inv}. Parmi ces personnes, 
nous notons la présence de 16 femmes et 49 hommes, soit 25\% de femmes. 
Compte-tenu du pourcentage de femmes dans la communauté, il y a une 
certaine forme de représentativité. Mais nous pourrions regretter que, 
dans certains cas, il y ait eu entre 3 et 5 personnes en conférence 
invitée, sans aucune femme (4 éditions concernées).

La présence en conférence invitée peut répondre à différentes logiques (y 
compris à celle du refus des personnes invitées) mais il apparaît 
nécessaire d'avoir une vigilance sur ce point, d'autant plus que la 
conférence EGC ne se limite pas à une seule invitation par édition.



\begin{figure}[h]\centering
\includegraphics[width=0.65\linewidth,keepaspectratio]{images/conf_inv.png}
\caption{Historique des conférences invitées : représentation 
sexuée.}\label{fig:conf_inv}
\end{figure}




\section{Focus environnement}\label{sec: environnement}

Au fil des années qui passent, l'environnement prend une place de plus en 
plus importante dans les débats questionnant l'organisation de nos 
sociétés. Parmi les préoccupations actuelles, on dénote sur un plan 
politique les émissions de gaz à effets de serre, principalement le CO2.

Dans le domaine de la recherche, une part des trajets à fortes émissions 
est en général attribuée aux conférences avec les déplacements nécessaires 
vers une zone géographique souvent éloignée de son lieu de travail commun. 
Différentes initiatives ou journées de 
réflexion\footnote{http://cedd-pes.com/nos-activites-2019/les-chercheurs-prennent-trop-lavion/} 
émergent sur ce sujet et se structurent petit à petit pour questionner la 
manière d'exercer une profession mais également la dimension écologique de 
l'informatique plus largement.


EcoInfo\footnote{https://ecoinfo.cnrs.fr/} est un groupement de services 
du CNRS qui a pour objectif général d’évaluer puis de réduire les impacts 
de l’informatique, notamment au sein de l’enseignement supérieur et de la 
recherche. Ce groupe a analysé ses propres émissions montrant la 
prépondérance des déplacements dans leurs émissions 
\citep{eco-info-carbone}.

Dans le cadre du défi EGC 2020, nous souhaitons pouvoir disposer d'une 
estimation en termes de distance et d'émission de carbone liées au 
déplacement pour la conférence. N'ayant pas accès aux participant·es de 
chaque édition, il s'agit pour nous d'aller vers le calcul d'un minimum en 
s'appuyant sur l'hypothèse que le déplacement a été effectué par la 
personne en première position sur le papier. Cette hypothèse de travail 
est nécessaire, même si elle peut être forte dans certaines situations où 
le papier est la production d'une collaboration entre des personnes 
d'organismes distants géographiquement. L'objectif visé n'est pas ici une 
estimation exacte, notamment car hormis les auteur/trices des papiers, 
d'autres personnes participent à la conférence, mais bien d'avoir une 
première idée des kilomètres parcourus et de ce que cela peut générer en 
termes d'émissions de carbone "au moins", en lien avec l'historique d'EGC.

En 2016, la publication de \cite{kergosien16egc} mentionnait une 
distribution des distances des auteur/trices au lieu de la conférence en 
ayant manuellement rempli leur base d'auteur/trices et ville du 
laboratoire de rattachement sans, par ailleurs, mettre en perspective la 
dimension écologique. Ici, nous nous proposons non seulement d'automatiser 
ce processus afin d'obtenir les distances parcourues par les personnes en 
première position sur les papiers pour l'ensemble du corpus de 
publications EGC depuis 2006, et d'aller par ailleurs vers une estimation 
totale minimale d'émissions de carbone en se basant sur des estimations 
kilométriques par type de transport (exprimées dans l'unité 
gCO2e/passager.km : grammes CO2 équivalent par passager et par kilomètre).



\subsection{Méthodologie d'extraction}\label{sec: methodologie}

Le jeu de données en entrée est constitué des informations DBLP sur les 
publications d'EGC et les PDF associés extraits de RNTI. La première étape 
consiste à extraire le texte de ces PDF pour obtenir les laboratoires de 
rattachement des auteur/trices. Afin de procéder à cette extraction, nous 
pouvons nous baser sur pdfminer. Cette démarche ne donne toutefois un 
résultat valide que pour 700 publications sur les 1390 analysées. Pour 
compenser ce manque, nous avons appliqué une méthode additionnelle sur les 
publications
problématiques. La méthode la plus fiable que nous ayons trouvée consiste 
à générer une image
de la première page de la publication avec la commande `convert' fournie 
par `image magick'.
Une fois cette image obtenue avec une résolution de 300dpi pour qu'elle 
soit suffisante, nous avons appliqué une reconnaissance de caractères avec 
l'outil `tesseract' rendu open source par HP il y a quelques années.
Sur nos tests, cet outil donne de bien meilleurs résultats que l'habituel 
GOCR fourni
par la communauté GNU, avec beaucoup moins d'erreurs dans les fichiers 
textes générés.

Pour s'assurer d'une bonne découpe des auteur/trices, nous effectuons une 
correspondance avec les informations fournies par les fichiers JSON de la 
base DBLP, ce qui permet de détecter les erreurs liées aux publications où 
les informations d'auteur/trices sont mal formattées. Ces dernières
ont été corrigées à la main (environ 300) pour pouvoir être repassées par 
les scripts automatiques. Les problèmes sont généralement liés aux 
virgules entre auteur/trices et caractères étoile (liant auteur/trice et 
laboratoire) mal reconnus.




\subsection{Gestion des distances}\label{sec: gestion}

Grâce à l'étape précédente, nous avons obtenu le laboratoire des 
auteur/trices et, pour ce qui
nous intéresse, la personne en première position sur l'article. \`{A} 
partir de cette information, nous pouvons utiliser les services de Google 
pour géolocaliser le laboratoire de manière automatique. Ceci nous donne 
une extraction viable de 916 publications sur les 1390 qui constitue notre 
base de travail pour construire les estimations minimales des métriques 
qui suivent. Les erreurs proviennent soit d'une mauvaise extraction des 
laboratoires depuis les papiers, soit d'un échec d'extraction de 
localisation reporté par Google. Ceci représente une perte de 34\% qui est 
relativement importante mais pourra être amélioré en ajoutant des fix 
automatiques sur les erreurs communes et récurrentes. Les taux de succès 
sont présentés dans le tableau \ref{fig: tableau}.

\begin{table}[h]
\centering \scriptsize{
		\begin{tabular}{|c|c|c|c||c|c|c|c|c|c|}
		\hline
Année &	Nb traités & Nb 
publiés & 	\% traitement & Année &	Nb traités & Nb publiés & 	\% 
traitement \\
\hline
2006 &	84	& 102 &	82\% & 2013 &	46	& 56 &	82\% \\
2007 &	77	& 92 &	84\% & 2014 &	67	& 87 &	77\% \\
2008 &	83	& 93 &	89\% & 2015 &	46	& 67 &	69\% \\ 
2009 &	74	& 82 &	90\% & 2016 &	61	& 84 &	73\%\\
2010 &	88	& 115 &	77\% & 2017 &	60	& 73 &	82\%\\
2011 &	84	& 100 &	84\% & 2018 &	53	& 71 &	75\%\\
2012 &	44	& 59 &	75\% & 2019 &	49	& 69 &	71\%\\
\hline
		\end{tabular}}
	\caption{Traitement des papiers pour le calcul de distances sur la 
période 2006-2019}\label{fig: tableau}
    
\end{table}



Nous pouvons sommer les distances associées aux publications par année et 
les comparer à la circonférence terrestre qui est d'environ 40 075 km 
\citep{earth}. Les résultats sont reportés dans le graphique 
\ref{tab:distance_glob}.
Toutes les distances supérieures à 800 km ont été vérifiées, nous obtenons 
donc un minimum considérant qu'il y a des publications non traitées, ainsi 
que l'hypothèse selon laquelle la personne en première position sur le 
papier fait le voyage.

\begin{figure}[h]\centering
\includegraphics[width=0.65\linewidth,keepaspectratio]{images/distance_glob.png}
\caption{Distance globale parcourue par édition, en km et en nombre de 
tours de la Terre.}\label{fig:distance_glob}
\end{figure}








\subsection{Empreinte carbone}\label{sec: ecologie}



Dans cette analyse, nous avons cherché une estimation minimale des 
émissions liées aux éditions EGC de 2006 à 2019 à des fins de 
sensibilisation. Pour parvenir à cette évaluation, nous avons été amené·es 
à simplifier la manière dont sont traités les trajets d'un point de vue de 
leur estimation selon les informations accessibles. Ainsi, deux modes de 
transport sont pris en compte : le train et l'avion. Nous considérerons 
(certes de façon arbitraire) le train si une gare est trouvée à moins de 
50 km du lieu de travail et du lieu de la conférence ; sinon, le moyen de 
transport considéré sera l'avion. Pour le train, le trajet considéré est 
celui entre les deux gares, le trajet pour aller à la gare étant ignoré. 
Pour le train, nous considérerons le type TGV (3,69 gCO2e/passager.km, 
source eco-info/ADEME). Les TER consommant davantage (8.91 
gCO2e/passager.km) seront ignorés en cherchant une valeur minimale des 
émissions. Concernant les avions, nous prendrons les bases de l'aviation 
civile\footnote{https://eco-calculateur.dta.aviation-civile.gouv.fr/} en 
considérant 2 types de trajets : ceux de plus de 1000km et ceux de moins 
de 1000km, les consommations à considérer étant différentes selon la 
catégorie. Nous prendrons donc pour ces deux catégories une moyenne de 
type Paris/Marseille (117 gCO2e/passager.km) et Paris/New York (82 
gCO2e/passager.km) respectivement. Cette estimation prendra en compte la 
distance entre le laboratoire et le lieu de la conférence, approximant le 
trajet aux aéroports par ces mêmes moyennes.


Pour les trajets en train, une estimation plus juste peut être obtenue en 
utilisant les données ouvertes et API publiques de la SNCF, avec la liste 
des gares et les API d'obtention des horaires avec les émissions 
associées. Cette méthode n'a pu être appliquée que sur 500 des 
publications. Ainsi, nous nous limitons ici à l'estimation avec les 
moyennes par moyen de transport. 



Les résultats donnent une estimation basse de 105 tonnes de CO2 équivalent 
émises pour le total des trajets considérés. Le maximum à Tunis donnant 15 
tonnes émises et les années les deux plus basses donnant respectivement 2 
tonnes en 2004 (Clermont Ferrand) et 3 tonnes en 2019 à Metz. Nous notons 
qu'une prise en compte des trajets réels effectués en train pour les 
extractions faites avec succès ne change pas significativement ces 
résultats, les nombres étant dominés par les trajets en avion. 

\section{Conclusion}\label{sec: conclusion}



Dans ce papier, nous avons choisi de porter un regard sur les données EGC 
au travers d'enjeux sociétaux d'actualité. Ceci constitue l'originalité de 
ce papier, même si les méthodes mobilisées ne sont pas nécessairement 
complexes. Les analyses ainsi réalisées sur les questions de la 
déconnexion dans le rapport au travail, de l'égalité femmes-hommes et 
d'enjeux environnementaux au travers de l'empreinte carbone fournissent 
des éléments éclairants.

La récupération de données autres que celles founies, de façon quasi 
automatique (dimension géographique) et manuelle (dimension égalité 
femmes-hommes) a nécessité un temps de travail important, restreignant les 
analyses menées qui pourraient donc être étendues.

Ainsi, en perspectives, des analyses pourraient être développées. Par 
exemple, sur le rapport au travail, une analyse centrée personne plutôt 
que sur le mail d'envoi serait pertinente, pour aller vers des éléments 
fins sur les profils de personnes (des adresses mails différentes peuvent 
correspondre à une seule et même personne : changement du nom de domaine 
de l'institution ou du nom en cas de mariage par exemple). En outre, nous 
avons centré notre analyse sur l'information des horaires d'envoi, une 
analyse au niveau des week-ends et des périodes de vacances serait 
intéressante par rapport à la définition du droit à la déconnexion. 
Concernant la dimension écologique, différents scénarios d'estimations 
pourraient être envisagés.



Pour finalement aller dans le sens d'une réponse au défi « EGC 2020: 20 
ans d’histoire pour quel avenir ? », les analyses fournies dans ce papier 
ne peuvent qu'inviter les chercheur·es de la communauté EGC et leurs 
responsables à se saisir des enjeux sociétaux qui traversent notre métier, 
notre posture professionnelle, à des niveaux à la fois individuel et 
collectif. Préserver son temps de sommeil, penser à (et agir pour) la 
place des femmes dans la communauté en termes de visibilité, faire le 
choix de l'écologie en matière de destination de conférences, autant de 
pistes qui nous amènent à une réflexivité sur nos manières d'exercer notre 
métier pour construire l'avenir d'EGC selon des choix conscients, et 
pourquoi pas engagés ?


\bibliographystyle{rnti}
\bibliography{biblio_EGC2020}




\Fr


\end{document}

